% =============================================================
% VTBench — Related Work (VIS Short Paper)
% =============================================================
% Usage: % =============================================================
% VTBench — Related Work (VIS Short Paper)
% =============================================================
% Usage: % =============================================================
% VTBench — Related Work (VIS Short Paper)
% =============================================================
% Usage: % =============================================================
% VTBench — Related Work (VIS Short Paper)
% =============================================================
% Usage: \input{related_work.tex} in your main .tex file
% Bibliography: \bibliography{references}
% =============================================================

\section{Related Work}
\label{sec:related}

Time series classification~(TSC) has traditionally relied on
numerical-domain representations and specialized architectures.
Early benchmarks such as the UCR archive~\cite{dau2019ucr} established a
standardized evaluation setting, while subsequent work explored deep neural
networks ranging from fully convolutional networks to residual architectures
for learning temporal patterns directly from raw
signals~\cite{fawaz2019deep}.
More recent studies extend this line to attention-based models and
self-supervised learning frameworks~\cite{foumani2024survey}.
In parallel, efficient alternatives such as ROCKET~\cite{dempster2020rocket}
and structured CNN designs~\cite{tang2022oscnn} demonstrate that strong
performance can be achieved without heavy model complexity.
These approaches operate directly on numerical sequences and serve as the
primary baselines for our study.

A complementary direction encodes time series as images, enabling the use of
computer vision models.
Early work introduced Gramian Angular Fields and Markov Transition Fields to
transform temporal signals into structured
images~\cite{wang2015imaging}.
Recurrence plots, originating from dynamical systems
theory~\cite{eckmann1987recurrence}, were later adopted for classification
with convolutional networks~\cite{hatami2018classification}.
Recent efforts further explore richer representations, including wavelet
transform-based scalograms~\cite{zhao2024cwt3dcnn} and hybrid encodings that
combine multiple image modalities~\cite{mariani2024fusion}.
In addition, rendering time series as chart images and leveraging pretrained
vision models has shown promising results~\cite{li2023vitst}.
More advanced encoding strategies, such as wavelet scattering
transforms~\cite{anden2014scattering}, path
signatures~\cite{kidger2021signatory}, and persistence images from
topological data analysis~\cite{seversky2016tda, umeda2017tda}, offer
theoretically grounded alternatives but have received limited comparative
evaluation.
However, these methods are typically studied in isolation and lack systematic
comparison across encoding strategies, model architectures~(including
pretrained backbones such as
EfficientNet~\cite{tan2019efficientnet}), and fusion configurations.

Interpretability has also emerged as an important concern in time series
analysis.
While numerical models provide limited transparency, image-based
representations offer a natural interface for visual inspection.
Gradient-based localization methods such as
Grad-CAM~\cite{selvaraju2017gradcam} have been widely used to highlight
discriminative regions in image models, providing insights into model
behavior.
Complementary work has investigated robustness through data augmentation
strategies in the time series
domain~\cite{iglesias2023augmentation}, and image-level post-processing
techniques (edge detection, contrast-limited adaptive histogram equalization)
can further enhance encoding discriminability, although their interaction
with visual encodings remains underexplored.

In this work, we provide a controlled empirical comparison that spans
numerical baselines, nine image encoding methods---including both classical
(GASF, MTF, RP) and advanced (scattering, signature, persistence)
representations---and four deep architectures, under a unified experimental
setting across 20 UCR datasets.
 in your main .tex file
% Bibliography: \bibliography{references}
% =============================================================

\section{Related Work}
\label{sec:related}

Time series classification~(TSC) has traditionally relied on
numerical-domain representations and specialized architectures.
Early benchmarks such as the UCR archive~\cite{dau2019ucr} established a
standardized evaluation setting, while subsequent work explored deep neural
networks ranging from fully convolutional networks to residual architectures
for learning temporal patterns directly from raw
signals~\cite{fawaz2019deep}.
More recent studies extend this line to attention-based models and
self-supervised learning frameworks~\cite{foumani2024survey}.
In parallel, efficient alternatives such as ROCKET~\cite{dempster2020rocket}
and structured CNN designs~\cite{tang2022oscnn} demonstrate that strong
performance can be achieved without heavy model complexity.
These approaches operate directly on numerical sequences and serve as the
primary baselines for our study.

A complementary direction encodes time series as images, enabling the use of
computer vision models.
Early work introduced Gramian Angular Fields and Markov Transition Fields to
transform temporal signals into structured
images~\cite{wang2015imaging}.
Recurrence plots, originating from dynamical systems
theory~\cite{eckmann1987recurrence}, were later adopted for classification
with convolutional networks~\cite{hatami2018classification}.
Recent efforts further explore richer representations, including wavelet
transform-based scalograms~\cite{zhao2024cwt3dcnn} and hybrid encodings that
combine multiple image modalities~\cite{mariani2024fusion}.
In addition, rendering time series as chart images and leveraging pretrained
vision models has shown promising results~\cite{li2023vitst}.
More advanced encoding strategies, such as wavelet scattering
transforms~\cite{anden2014scattering}, path
signatures~\cite{kidger2021signatory}, and persistence images from
topological data analysis~\cite{seversky2016tda, umeda2017tda}, offer
theoretically grounded alternatives but have received limited comparative
evaluation.
However, these methods are typically studied in isolation and lack systematic
comparison across encoding strategies, model architectures~(including
pretrained backbones such as
EfficientNet~\cite{tan2019efficientnet}), and fusion configurations.

Interpretability has also emerged as an important concern in time series
analysis.
While numerical models provide limited transparency, image-based
representations offer a natural interface for visual inspection.
Gradient-based localization methods such as
Grad-CAM~\cite{selvaraju2017gradcam} have been widely used to highlight
discriminative regions in image models, providing insights into model
behavior.
Complementary work has investigated robustness through data augmentation
strategies in the time series
domain~\cite{iglesias2023augmentation}, and image-level post-processing
techniques (edge detection, contrast-limited adaptive histogram equalization)
can further enhance encoding discriminability, although their interaction
with visual encodings remains underexplored.

In this work, we provide a controlled empirical comparison that spans
numerical baselines, nine image encoding methods---including both classical
(GASF, MTF, RP) and advanced (scattering, signature, persistence)
representations---and four deep architectures, under a unified experimental
setting across 20 UCR datasets.
 in your main .tex file
% Bibliography: \bibliography{references}
% =============================================================

\section{Related Work}
\label{sec:related}

Time series classification~(TSC) has traditionally relied on
numerical-domain representations and specialized architectures.
Early benchmarks such as the UCR archive~\cite{dau2019ucr} established a
standardized evaluation setting, while subsequent work explored deep neural
networks ranging from fully convolutional networks to residual architectures
for learning temporal patterns directly from raw
signals~\cite{fawaz2019deep}.
More recent studies extend this line to attention-based models and
self-supervised learning frameworks~\cite{foumani2024survey}.
In parallel, efficient alternatives such as ROCKET~\cite{dempster2020rocket}
and structured CNN designs~\cite{tang2022oscnn} demonstrate that strong
performance can be achieved without heavy model complexity.
These approaches operate directly on numerical sequences and serve as the
primary baselines for our study.

A complementary direction encodes time series as images, enabling the use of
computer vision models.
Early work introduced Gramian Angular Fields and Markov Transition Fields to
transform temporal signals into structured
images~\cite{wang2015imaging}.
Recurrence plots, originating from dynamical systems
theory~\cite{eckmann1987recurrence}, were later adopted for classification
with convolutional networks~\cite{hatami2018classification}.
Recent efforts further explore richer representations, including wavelet
transform-based scalograms~\cite{zhao2024cwt3dcnn} and hybrid encodings that
combine multiple image modalities~\cite{mariani2024fusion}.
In addition, rendering time series as chart images and leveraging pretrained
vision models has shown promising results~\cite{li2023vitst}.
More advanced encoding strategies, such as wavelet scattering
transforms~\cite{anden2014scattering}, path
signatures~\cite{kidger2021signatory}, and persistence images from
topological data analysis~\cite{seversky2016tda, umeda2017tda}, offer
theoretically grounded alternatives but have received limited comparative
evaluation.
However, these methods are typically studied in isolation and lack systematic
comparison across encoding strategies, model architectures~(including
pretrained backbones such as
EfficientNet~\cite{tan2019efficientnet}), and fusion configurations.

Interpretability has also emerged as an important concern in time series
analysis.
While numerical models provide limited transparency, image-based
representations offer a natural interface for visual inspection.
Gradient-based localization methods such as
Grad-CAM~\cite{selvaraju2017gradcam} have been widely used to highlight
discriminative regions in image models, providing insights into model
behavior.
Complementary work has investigated robustness through data augmentation
strategies in the time series
domain~\cite{iglesias2023augmentation}, and image-level post-processing
techniques (edge detection, contrast-limited adaptive histogram equalization)
can further enhance encoding discriminability, although their interaction
with visual encodings remains underexplored.

In this work, we provide a controlled empirical comparison that spans
numerical baselines, nine image encoding methods---including both classical
(GASF, MTF, RP) and advanced (scattering, signature, persistence)
representations---and four deep architectures, under a unified experimental
setting across 20 UCR datasets.
 in your main .tex file
% Bibliography: \bibliography{references}
% =============================================================

\section{Related Work}
\label{sec:related}

Time series classification~(TSC) has traditionally relied on
numerical-domain representations and specialized architectures.
Early benchmarks such as the UCR archive~\cite{dau2019ucr} established a
standardized evaluation setting, while subsequent work explored deep neural
networks ranging from fully convolutional networks to residual architectures
for learning temporal patterns directly from raw
signals~\cite{fawaz2019deep}.
More recent studies extend this line to attention-based models and
self-supervised learning frameworks~\cite{foumani2024survey}.
In parallel, efficient alternatives such as ROCKET~\cite{dempster2020rocket}
and structured CNN designs~\cite{tang2022oscnn} demonstrate that strong
performance can be achieved without heavy model complexity.
These approaches operate directly on numerical sequences and serve as the
primary baselines for our study.

A complementary direction encodes time series as images, enabling the use of
computer vision models.
Early work introduced Gramian Angular Fields and Markov Transition Fields to
transform temporal signals into structured
images~\cite{wang2015imaging}.
Recurrence plots, originating from dynamical systems
theory~\cite{eckmann1987recurrence}, were later adopted for classification
with convolutional networks~\cite{hatami2018classification}.
Recent efforts further explore richer representations, including wavelet
transform-based scalograms~\cite{zhao2024cwt3dcnn} and hybrid encodings that
combine multiple image modalities~\cite{mariani2024fusion}.
In addition, rendering time series as chart images and leveraging pretrained
vision models has shown promising results~\cite{li2023vitst}.
More advanced encoding strategies, such as wavelet scattering
transforms~\cite{anden2014scattering}, path
signatures~\cite{kidger2021signatory}, and persistence images from
topological data analysis~\cite{seversky2016tda, umeda2017tda}, offer
theoretically grounded alternatives but have received limited comparative
evaluation.
However, these methods are typically studied in isolation and lack systematic
comparison across encoding strategies, model architectures~(including
pretrained backbones such as
EfficientNet~\cite{tan2019efficientnet}), and fusion configurations.

Interpretability has also emerged as an important concern in time series
analysis.
While numerical models provide limited transparency, image-based
representations offer a natural interface for visual inspection.
Gradient-based localization methods such as
Grad-CAM~\cite{selvaraju2017gradcam} have been widely used to highlight
discriminative regions in image models, providing insights into model
behavior.
Complementary work has investigated robustness through data augmentation
strategies in the time series
domain~\cite{iglesias2023augmentation}, and image-level post-processing
techniques (edge detection, contrast-limited adaptive histogram equalization)
can further enhance encoding discriminability, although their interaction
with visual encodings remains underexplored.

In this work, we provide a controlled empirical comparison that spans
numerical baselines, nine image encoding methods---including both classical
(GASF, MTF, RP) and advanced (scattering, signature, persistence)
representations---and four deep architectures, under a unified experimental
setting across 20 UCR datasets.
